\documentclass[11pt]{article}
\usepackage{amsmath,amssymb}
\usepackage[margin=0.75in]{geometry}
\usepackage{enumitem}

\setlist{leftmargin=*,nosep}
\setlength{\parskip}{2pt}

\begin{document}
\begin{center}
MAT067 \quad University of California, Davis \quad Winter 2010

\Large\textbf{Midterm Exam}\\[4pt]
\normalsize (Wednesday, February 10)
\end{center}

\small
\section*{Problem 1 (15 points)}
\begin{enumerate}[label=\alph*)]
  \item Suppose we are given a complex number $z = x + iy$ where $x, y \in \mathbb{R}$. Write the definitions of the complex conjugate $\overline{z}$ and the modulus $|z|$.
  \item Using only part~(a), prove that for nonzero $z$ the number $w = \dfrac{\overline{z}}{|z|^{2}}$ is an inverse of $z$.
  \item Prove that $e^{i\theta_{1}} e^{i\theta_{2}} = e^{i(\theta_{1} + \theta_{2})}$ using only the polar form definition of $e^{i\theta}$.
  \item Find all solutions over $\mathbf{C}$ to $z^{3} = -1$.
\end{enumerate}

\section*{Problem 2 (15 points)}
Let $V$ be a vector space over $\mathbf{F}$. (Parts d--e continue on page 3.)
\begin{enumerate}[label=\alph*)]
  \item State the three conditions for a subset $U \subseteq V$ to be a subspace of $V$.
  \item Let $U_{1}$ and $U_{2}$ be subspaces of $V$. Which of the above conditions prevents $U_{1} \cup U_{2} = \{v \in V \mid v \in U_{1} \text{ or } v \in U_{2}\}$ from being a subspace of $V$?
  \item Let $U_{1}$ and $U_{2}$ be subspaces of $V$. Show that if $V = U_{1} \oplus U_{2}$ (direct sum) then $U_{1} \cap U_{2} = \{0\}$.
  \item Let $V = \mathbf{C}^{4}$ and $U_{1} = \{(a,b,c,d) \in V : d = 0\}$. Show that $U_{1}$ is a subspace of $V$.
  \item Show that $V$ is not the direct sum of $U_{1}$ and $U_{2}$ where $U_{2} = \{(a,b,c,d) \in V : a + 2b + 3c + 4d = 0\}$.
\end{enumerate}

\section*{Problem 3 (15 points)}
Suppose $V$ and $W$ are vector spaces over $\mathbf{F}$ and $T: V \to W$ is linear.
\begin{enumerate}[label=\alph*)]
  \item Write the definitions of $\operatorname{null}(T)$ and $\operatorname{range}(T)$.
  \item Suppose we have found a subspace $U$ such that $V = \operatorname{null}(T) \oplus U$. Show that $\operatorname{range}(T) = \{T(u) : u \in U\}$.
  \item In the setting of part~(b), show that $\dim(\operatorname{range}(T)) = \dim U$.
  \item Using the formula from part~(b) and/or (c), or other methods, find $\dim(\operatorname{null}(T))$ for the map $T : \mathbf{R}^{4} \to \mathbf{R}^{3}$ given by
  \[
    T:\;
    \begin{pmatrix} a \\ b \\ c \\ d \end{pmatrix}
    \longmapsto
    \begin{pmatrix}
      1 & 1 & 4 & 3 \\
      2 & 1 & 3 & 2 \\
      3 & 1 & 2 & 1
    \end{pmatrix}
    \begin{pmatrix} a \\ b \\ c \\ d \end{pmatrix}.
  \]
\end{enumerate}
\medskip
For part~(d) you may use any linear relations among the listed $v_{i}$ vectors that you can justify.

\end{document}
